% dynamics/dynamic_balance-DE.tex
\chapter{Dynamisches Gleichgewicht als qualifiziertes Residuum}
\label{chap:dynamic-balance}

\section*{Eine Rechnung, sechs verschiedene Rollen}

Betrachtet wird der reduzierte Punktfolgetest \(M_{\mathrm{red}}\) im lokalen
Gleis-Frame \(F_{\mathrm{track}}\). Positive Krümmung und positive Querwirkung
verwenden dieselbe erklärte Seite; ein positiver Überhöhungswinkel \(\phi\)
kompensiert diese Wirkung. Der Test setzt starre Durchfahrt, vorgegebenes
\(v(t)\), geringe lokale Ausdehnung und einheitlichen Schwerebetrag \(g\)
voraus. Er lässt Fahrzeugmasse und -trägheit, Federung, Rad--Schiene-Kontakt,
Gleisnachgiebigkeit, Unregelmäßigkeiten, Aerodynamik und Regelung aus.

\begin{principle}[Prinzip des dynamischen Gleichgewichts]
\label{princ:dynamic-balance}
Dynamisches Gleichgewicht ist eine nach Modell, Frame, Vorzeichen und Zweck
qualifizierte Relation.
\end{principle}

\section{Gleichgewichtsrelation}

\begin{definition}[Gleichgewichtsrelation]
\label{def:balance-relation}
Für \(M_{\mathrm{red}}\) lautet das Gleichgewichtsresiduum
\[
r_{\mathrm{bal}}(t)
=v(t)^2\kappa(s(t))-g\sin\phi(s(t)).
\]
Die reduzierte Gleichgewichtsbedingung ist \(r_{\mathrm{bal}}=0\).
\end{definition}

\[
\boxed{\begin{gathered}
\text{verschwindendes reduziertes}\\
\text{Gleichgewichtsresiduum}
\end{gathered}}
\not\Rightarrow
\boxed{\begin{gathered}
\text{verschwindende}\\
\text{Wagenkastenbeschleunigung}
\end{gathered}}
\]

Die Gleichheit hebt zwei Terme in diesem einen Frame und Modell auf. Sie
begründet weder Mehrkörpergleichgewicht noch vorhergesagte Fahrzeugantwort,
Messantwort, Fahrgastwahrnehmung, Sicherheit oder Komfort.

\begin{definition}[Effektive Querbeschleunigung]
\label{def:dynamics-effective-lateral-acceleration}
Die historische Schreibweise \(a_{\mathrm{eff}}\) bezeichnet in diesem
reduzierten Test \(r_{\mathrm{bal}}\). Die Benennung als Beschleunigung
erweitert ihren Evidenzgehalt nicht.
\end{definition}

\begin{figure}[htbp]
\centering
% figures/dynamics/dynamic_balance_relation-DE.tex
\begin{tikzpicture}[
  node distance=1.6cm,
  box/.style={draw, rounded corners, minimum width=3.8cm, minimum height=0.9cm, align=center},
  smallbox/.style={draw, rounded corners, minimum width=3.4cm, minimum height=0.9cm, align=center},
  arrow/.style={->, thick}
]
\node[box] (curv) {Krümmung\\$\kappa(s)$};
\node[box, right=of curv] (acc) {Querbeschleunigung\\$v^2\kappa(s)$};
\node[box, below=of acc] (cant) {Überhöhungskompensation\\$g\sin\phi(s)$};
\node[box, right=of acc] (balance) {Gleichgewichtsbeziehung\\$v^2\kappa(s)=g\sin\phi(s)$};
\node[smallbox, below=of balance] (imbalance) {Ungleichgewicht\\$a_{\mathrm{eff}}=v^2\kappa-g\sin\phi$};
\draw[arrow] (curv) -- (acc);
\draw[arrow] (acc) -- (balance);
\draw[arrow] (cant) -- (balance);
\draw[arrow] (balance) -- (imbalance);
\end{tikzpicture}

\caption{Das reduzierte Gleis-Frame-Residuum; ausgelassene Antwortphysik bleibt außerhalb der Relation.}
\label{fig:dynamic-balance-relation}
\end{figure}

\section{Überhöhungsfehlbetrag, -überschuss und Verlauf}

\begin{definition}[Überhöhungsfehlbetrag]
\label{def:cant-deficiency}
Unter den erklärten Vorzeichen und \(M_{\mathrm{red}}\) ist
\(r_{\mathrm{bal}}>0\) ein reduzierter Überhöhungsfehlbetrag. Eine
standardspezifische Größe benötigt weiterhin eigene Definition und
Anwendbarkeit.
\end{definition}

\begin{definition}[Überhöhungsüberschuss]
\label{def:excess-cant}
Unter denselben Qualifikationen ist \(r_{\mathrm{bal}}<0\) ein reduzierter
Überhöhungsüberschuss.
\end{definition}

\begin{definition}[Gleichgewichtstrajektorie]
\label{def:dynamics-balance-trajectory}
Die Gleichgewichtstrajektorie ist \(r_{\mathrm{bal}}(t)\), oder nur dort
\(r_{\mathrm{bal}}(s)\), wo die Geschwindigkeit als Stationsfunktion
vorgegeben wurde.
\end{definition}

Bei veränderlicher Geschwindigkeit gilt
\[
\dot r_{\mathrm{bal}}
=2v\dot v\kappa+v^3\kappa'-gv\cos\phi\,\phi',
\]
unter denselben Annahmen. Die Ableitung ergänzt die Expositionsrate, aber
weiterhin keine Fahrzeugantwort.

\section{Die Rollen nicht zusammenziehen}

\begin{center}
\small
\begin{tabular}{p{0.19\linewidth}p{0.33\linewidth}p{0.37\linewidth}}
\toprule
Rolle & Erzeugt & Begründet nicht\\
\midrule
Residuum & \(r_{\mathrm{bal}}\) aus \(M_{\mathrm{red}}\) & Wagenkastenantwort\\
Vorhersage & \(x_M^{\mathrm{pred}}(t)\) aus identifiziertem \(M\) & Beobachtung\\
Beobachtung & quellen- und unsicherheitsqualifizierte Schätzung & physische Wahrheit oder Ursache\\
Vergleich & Differenz bei gleicher Identität, Zeit und gleichem Frame & Akzeptabilität\\
Bewertung & Kriterienergebnis für einen Zweck & Auswahl oder Autorität\\
Engineering Decision & dokumentierte Auswahl oder Ablehnung & universelle Gültigkeit\\
\bottomrule
\end{tabular}
\end{center}

\[
\boxed{\text{vorhergesagte Antwort}
\not\Rightarrow \text{akzeptabler ingenieurtechnischer Zustand}}
\]

CAD-Kurve und Vermessungspolylinie machen die Trennung konkret. Beide können
zu einer vergleichbaren Gleis-Frame-Krümmungsschätzung realisiert werden;
eine Abweichung ist jedoch nicht automatisch Ursache einer gemessenen
Wagenkastenbeschleunigung. Diese Kausalaussage benötigt ein identifiziertes
Antwortmodell, zugeordnete Beobachtungen und qualifizierte Evidenz. Mit genau
dieser Zurückhaltung wird das folgende Vienna-Beispiel gelesen.
